\section{The exact \texorpdfstring{$\mathbb Z_4$}{Z4} tower}
\label{sec:tower}

Let
\[
 \A_1\xrightarrow{\iota_1}\A_2\xrightarrow{\iota_2}\cdots
\]
be a unital inductive system of finite matrix algebras.  Suppose
$u_N\in\A_N$ is unitary,
\[
 u_N^4=\id,\qquad \iota_N(u_N)=u_{N+1},
\tag{4.1}\label{eq:compatible-clock}
\]
and let $\alpha_N=\operatorname{Ad}(u_N)$.
Write $v_N$ for the canonical crossed-product unitary:
\[
 v_Nav_N^*=\alpha_N(a),\qquad v_N^4=\id.
\]

\begin{theorem}[Naturality, fusion, and expectations]
\label{thm:crossed}
Under \eqref{eq:compatible-clock}:
\begin{enumerate}[label=\textup{(\alph*)}]
\item There is a natural $*$-isomorphism
\[
 \Phi_N:\A_N\rtimes_{\alpha_N}\mathbb Z_4
 \longrightarrow
 \A_N\otimes C^*(\mathbb Z_4),
\qquad
 \Phi_N(av_N^g)=au_N^g\otimes w^g,
\tag{4.2}\label{eq:phi}
\]
where $w$ is the canonical generator of $C^*(\mathbb Z_4)$.
\item If $\A=\varinjlim\A_N$ and $\alpha=\varinjlim\alpha_N$, then
\[
 \varinjlim(\A_N\rtimes_{\alpha_N}\mathbb Z_4)
 \cong\A\rtimes_\alpha\mathbb Z_4
 \cong\A\otimes C^*(\mathbb Z_4).
\tag{4.3}\label{eq:limit-crossed}
\]
\item The homogeneous subspaces
$\mathcal B_g=\overline{\A v^g}$ obey
\[
 \mathcal B_g\mathcal B_h\subset\mathcal B_{g+h},
\qquad
\mathcal B_g^*=\mathcal B_{-g},
\qquad
v^4=\id\in\A.
\tag{4.4}\label{eq:fusion}
\]
Thus a degree-one symbol $\Psi_\lambda:=v$ has
$\Psi_\lambda^4$ in the neutral algebra, and
$\Psi_\lambda^2$ is the unique order-two grade.  With the frozen label
$[\lambda]^2=[v]$, this is exactly the abstract $\mathbb Z_4$ fusion
law.
\item The coefficient map
\[
 F\left(\sum_{g=0}^3a_gv^g\right)=a_0
\tag{4.5}\label{eq:coefficient-expectation}
\]
extends to a faithful conditional expectation
$F:\A\rtimes_\alpha\mathbb Z_4\to\A$.  It is the unique conditional
expectation invariant under the dual $\widehat{\mathbb Z_4}$ action.
\item If $\tau_N$ is a compatible faithful trace on the tower, then
\[
 E_N(a)=\frac14\sum_{g=0}^3\alpha_N^g(a)
\tag{4.6}\label{eq:average-expectation}
\]
is the unique $\tau_N$-preserving conditional expectation
$\A_N\to\A_N^{\mathbb Z_4}$.  The $E_N$ are natural and induce the
corresponding trace-preserving expectation on the inductive limit.
\end{enumerate}
\end{theorem}

\begin{proof}
Map $a\in\A_N$ to $a\otimes\id$ and $v_N$ to $u_N\otimes w$.
The defining relations are preserved:
\[
\begin{aligned}
(u_N\otimes w)(a\otimes\id)(u_N^*\otimes w^*)
 &=u_Nau_N^*\otimes\id=\alpha_N(a)\otimes\id,\\
(u_N\otimes w)^4&=u_N^4\otimes w^4=\id.
\end{aligned}
\]
The universal property of the crossed product therefore gives a
$*$-homomorphism $\Phi_N$.  It is onto because its range contains
$a\otimes\id$ and also
\[
 \Phi_N(u_N^*v_N)=\id\otimes w.
\]
An inverse is given on generators by
\[
 a\otimes\id\longmapsto a,\qquad
 \id\otimes w\longmapsto u_N^*v_N.
\]
Indeed, $u_N^*v_N$ commutes with $\A_N$:
\[
 (u_N^*v_N)a
 =u_N^*\alpha_N(a)v_N
 =u_N^*u_Nau_N^*v_N
 =a(u_N^*v_N),
\]
and $(u_N^*v_N)^4=\id$ because $v_Nu_Nv_N^*=u_N$.
Thus \eqref{eq:phi} is an isomorphism.

For naturality, let
$\widetilde\iota_N:\A_N\rtimes\mathbb Z_4
\to\A_{N+1}\rtimes\mathbb Z_4$
send $a$ to $\iota_N(a)$ and $v_N$ to $v_{N+1}$.  On a spanning
monomial,
\[
\begin{aligned}
\Phi_{N+1}\widetilde\iota_N(av_N^g)
 &=\iota_N(a)u_{N+1}^g\otimes w^g\\
 &=\iota_N(au_N^g)\otimes w^g\\
 &=(\iota_N\otimes\id)\Phi_N(av_N^g).
\end{aligned}
\tag{4.7}\label{eq:naturality}
\]
So the diagrams commute exactly.  Taking inductive limits in
\eqref{eq:naturality} gives the first isomorphism in
\eqref{eq:limit-crossed}; applying the same generator argument to the
limit unitary $u=\varinjlim u_N$ gives the second.

For $a,b\in\A$,
\[
 (av^g)(bv^h)=a\alpha^g(b)v^{g+h}\in\A v^{g+h},
\]
and
\[
 (av^g)^*=v^{-g}a^*=\alpha^{-g}(a^*)v^{-g}.
\]
These identities prove \eqref{eq:fusion}, including $v^4=\id$.
The order-two identification is precisely the frozen lattice equality
$[\lambda]^2=[v]$; no net identification is inferred from the notation.

Under $\Phi$, the map $F$ is
\[
 \id_\A\otimes\tau_{\mathbb Z_4},
\qquad
\tau_{\mathbb Z_4}(w^g)=\delta_{g0}.
\]
The group trace is positive and faithful on the finite-dimensional
algebra $C^*(\mathbb Z_4)$, hence $F$ is positive, unital, faithful,
idempotent, fixes $\A$, and satisfies
$F(aXb)=aF(X)b$ for $a,b\in\A$.  Thus it is a faithful conditional
expectation.

Let $\widehat\alpha_\chi$ be the dual action,
$\widehat\alpha_\chi(a)=a$ and
$\widehat\alpha_\chi(v)=\chi(1)v$.  If $G$ is a dual-invariant
conditional expectation onto $\A$, then for $g\ne0$ choose a character
$\chi$ with $\chi(g)\ne1$.  For every $a\in\A$,
\[
 G(av^g)
 =G(\widehat\alpha_\chi(av^g))
 =\chi(g)G(av^g),
\]
so $G(av^g)=0$.  Also $G(a)=a$.  Hence $G=F$ on the algebraic crossed
product, a dense subalgebra, and continuity gives $G=F$ everywhere.

Finally, \eqref{eq:average-expectation} is plainly unital and fixes
$\A_N^{\mathbb Z_4}$.  It is positive because it is an average of
$*$-automorphisms, idempotent because
\[
 E_N(E_N(a))
 =\frac1{16}\sum_{g,h=0}^3\alpha_N^{g+h}(a)
 =\frac14\sum_{\ell=0}^3\alpha_N^\ell(a)=E_N(a),
\]
and bimodular because, for fixed $x,y$,
\[
 E_N(xay)=\frac14\sum_gx\alpha_N^g(a)y=xE_N(a)y.
\]
It preserves $\tau_N$ because traces are invariant under inner
automorphisms.

To prove uniqueness, equip $\A_N$ with
$\ip{x}{y}_{2,\tau_N}=\tau_N(x^*y)$.  For $b$ fixed by the action,
\[
\begin{aligned}
\ip{b}{a-E_N(a)}_{2,\tau_N}
&=\tau_N(b^*a)
 -\frac14\sum_g\tau_N(b^*\alpha_N^g(a))\\
&=\tau_N(b^*a)
 -\frac14\sum_g\tau_N(\alpha_N^{-g}(b^*)a)\\
&=\tau_N(b^*a)-\tau_N(b^*a)=0.
\end{aligned}
\]
Thus $E_N(a)$ is the orthogonal projection of $a$ onto the fixed
subspace.  If $G_N$ is another $\tau_N$-preserving conditional
expectation, then for fixed $b$ the bimodule property gives
\[
\ip{b}{a-G_N(a)}_{2,\tau_N}
=\tau_N(b^*a)-\tau_N(G_N(b^*a))
=\tau_N(b^*a)-\tau_N(b^*G_N(a))=0.
\]
So $G_N(a)$ is the same orthogonal projection, and $G_N=E_N$.
Compatibility follows by substituting
\eqref{eq:compatible-clock} into \eqref{eq:average-expectation}.
\end{proof}

\begin{remark}[Why the qualifier on uniqueness is necessary]
For an inner action,
$\A\rtimes\mathbb Z_4\cong\A\otimes C^*(\mathbb Z_4)$.
An arbitrary state on the second tensor factor defines a conditional
expectation onto $\A$, so a bare claim of uniqueness is false.
The coefficient expectation is unique after imposing dual invariance
(equivalently, the canonical group trace).  Separately, the averaging
expectation $\A\to\A^{\mathbb Z_4}$ is unique after imposing preservation
of the specified trace.  These are the two precise uniqueness statements
supported by the finite data.
\end{remark}

\section{The quarter-twist logarithm and global obstruction}
\label{sec:global}

\subsection{One chiral edge}

Let $\mathfrak h=L^2([0,1])$.  For $n,m\in\Z$ and
$0<\alpha<1$, define
\[
 e_n(x)=e^{2\pi inx},\qquad
 f_m^{(\alpha)}(x)=e^{2\pi i(m+\alpha)x}.
\]
Both families are orthonormal bases.  Let $P_0$ be the projection onto
the closed span of $\{e_n:n\le0\}$ and $P_\alpha$ the projection onto
the closed span of $\{f_m^{(\alpha)}:m\le0\}$.  The shift
$\alpha=k/4$ is exactly the momentum displacement
\[
 \Delta p=\frac{2\pi k}{4N_x}
\]
relative to the spacing $2\pi/N_x$.

\begin{lemma}[Exact overlap]\label{lem:overlap}
For all $m,n\in\Z$,
\[
 \left|\ip{e_n}{f_m^{(\alpha)}}\right|^2
 =\frac{\sin^2(\pi\alpha)}
 {\pi^2(m-n+\alpha)^2}.
\tag{5.1}\label{eq:overlap}
\]
\end{lemma}

\begin{proof}
Direct integration gives
\[
\begin{aligned}
\ip{e_n}{f_m^{(\alpha)}}
&=\int_0^1e^{2\pi i(m-n+\alpha)x}\,dx\\
&=\frac{e^{2\pi i(m-n+\alpha)}-1}
 {2\pi i(m-n+\alpha)}\\
&=e^{\pi i(m-n+\alpha)}
\frac{\sin(\pi(m-n+\alpha))}
 {\pi(m-n+\alpha)}.
\end{aligned}
\]
Since $m-n$ is an integer,
$|\sin(\pi(m-n+\alpha))|=|\sin(\pi\alpha)|$.
Taking the squared modulus proves \eqref{eq:overlap}.
\end{proof}

\begin{theorem}[One-edge Hilbert--Schmidt asymptotic]
\label{thm:oneedge}
Let $S_N(\alpha)$ be the contribution to
$\norm{P_0-P_\alpha}_2^2$ obtained by restricting both occupied and
unoccupied indices to absolute value at most $N$.  Then
\[
 S_N(\alpha)
 =\frac{2\sin^2(\pi\alpha)}{\pi^2}\log N+O_\alpha(1).
\tag{5.2}\label{eq:one-edge-log}
\]
In particular, $P_0-P_\alpha$ is not Hilbert--Schmidt.
\end{theorem}

\begin{proof}
For any two orthogonal projections $P,Q$ for which the terms are
defined,
\[
 (P-Q)^2=P(\id-Q)P+(\id-P)Q(\id-P).
\]
Taking the trace in orthonormal bases gives
\[
 \norm{P-Q}_2^2
 =\Tr(P(\id-Q))+\Tr(Q(\id-P)).
\tag{5.3}\label{eq:projection-distance}
\]
For the first term in \eqref{eq:projection-distance}, occupied
$e_n$ have $n\le0$ and unoccupied $f_m^{(\alpha)}$ have $m\ge1$.
For the second, occupied $f_m^{(\alpha)}$ have $m\le0$ and
unoccupied $e_n$ have $n\ge1$.  By Lemma~\ref{lem:overlap}, the
symmetrically truncated sum is
\[
 S_N(\alpha)
 =\frac{\sin^2(\pi\alpha)}{\pi^2}
 \sum_{d=1}^{2N}\mu_N(d)
 \left(\frac1{(d+\alpha)^2}+\frac1{(d-\alpha)^2}\right),
\tag{5.4}\label{eq:finite-sum}
\]
where
\[
 \mu_N(d)=
 \#\{(m,n):1\le m\le N,\ -N\le n\le0,\ m-n=d\}.
\]
Counting the integer points on each diagonal gives
\[
 \mu_N(d)=
 \begin{cases}
 d,&1\le d\le N,\\
 2N-d+1,&N<d\le2N.
 \end{cases}
\tag{5.5}\label{eq:multiplicity}
\]
The possible one-unit convention change at $d=N+1$ affects only a
bounded term and not the asymptotic; \eqref{eq:multiplicity} corresponds
to the displayed index ranges.

For $d\ge1$,
\[
\begin{aligned}
d\left(\frac1{(d+\alpha)^2}+\frac1{(d-\alpha)^2}\right)
&=\frac1d\left(
\frac1{(1+\alpha/d)^2}
+\frac1{(1-\alpha/d)^2}\right)\\
&=\frac2d+O_\alpha(d^{-3}).
\end{aligned}
\tag{5.6}\label{eq:summand-expand}
\]
The absence of a $d^{-2}$ term follows because the terms odd in
$\alpha/d$ cancel.  Therefore the part $1\le d\le N$ of
\eqref{eq:finite-sum} equals
\[
 2\sum_{d=1}^N\frac1d+O_\alpha\left(\sum_{d=1}^\infty d^{-3}\right)
 =2\log N+O_\alpha(1).
\tag{5.7}\label{eq:harmonic}
\]
For the tail $N<d\le2N$, use $d-\alpha\ge d/2$ to obtain
\[
\begin{aligned}
\sum_{d=N+1}^{2N}\mu_N(d)
\left(\frac1{(d+\alpha)^2}+\frac1{(d-\alpha)^2}\right)
&\le
\sum_{d=N+1}^{2N}(2N-d+1)\frac8{d^2}\\
&\le \frac8{N^2}\sum_{r=1}^{N}r
\le 4+\frac4N.
\end{aligned}
\tag{5.8}\label{eq:tail}
\]
Thus the tail is uniformly bounded.  Substitution of
\eqref{eq:harmonic} and \eqref{eq:tail} into \eqref{eq:finite-sum}
proves \eqref{eq:one-edge-log}.  Since the nonnegative truncations
diverge, the full sum in \eqref{eq:projection-distance} is infinite,
so $P_0-P_\alpha$ is not Hilbert--Schmidt.
\end{proof}

